\chapter*{第三部分导读：把机器学习当作科学推断工具}
\addcontentsline{toc}{chapter}{第三部分导读：把机器学习当作科学推断工具}

Part III 进入机器学习主线。这里最容易出现的误解，是把机器学习看成一组模型名称：线性回归、逻辑回归、随机森林、PCA、聚类、异常检测。模型当然要学，但本书更关心的是另一件事：如何把模型放进科学推断流程里，让它回答一个明确问题，并且知道它什么时候不该被相信。

从 Part II 过来之后，你已经知道数据带着单位、误差、采样、坐标、选择效应和处理记录。机器学习不会自动抹掉这些背景。相反，它会把这些背景压缩进特征、标签和样本划分里。如果这些入口没有想清楚，模型越强，越可能把问题藏得更深。

\section*{本部分的主线}

Part III 的章节顺序是一条从问题定义到可信解释的链条：

\begin{enumerate}
    \item 先区分 AI、机器学习和科学推断的关系；
    \item 再建立最小机器学习 workflow：任务定义、切分、baseline、训练、评估；
    \item 分别学习回归与分类这两类最常见监督任务；
    \item 用模型评估、诊断、特征工程和超参数选择控制可靠性；
    \item 进入无监督学习与异常检测，理解没有标签时如何探索结构；
    \item 最后讨论不确定度、解释性和科学可信度。
\end{enumerate}

这条主线有一个重要原则：\textbf{模型选择晚于问题定义，分数解释晚于数据理解。}

\section*{正文和训练材料如何配合}

Part III 的正文负责让你\textbf{看懂}机器学习方法在科学推断中的位置：模型在优化什么、baseline 为什么必要、评估为什么容易自欺、错误样本怎样回到科学问题、解释性和不确定度为什么不能事后才补。它不是模型名录，而是一套判断模型结果是否可信的语言。

配套训练材料则负责让你\textbf{练会}模型记录。每章都会产生一种小型 record：问题定义、workflow、regression、classification、evaluation、preprocessing、model selection、unsupervised、anomaly review 或 trust statement。后续案例章节不应重新发明这些记录，而应把它们组合成完整项目包。

\section*{贯穿 Part III 的三道防线}

为了避免把机器学习学成“调用库函数”，本部分始终保留三道防线。

第一道防线是 \textbf{baseline}。任何复杂模型之前，都应先有一个简单、可解释、容易复现的基线结果。没有 baseline，模型分数就很难判断是好是坏。

第二道防线是 \textbf{held-out evaluation}。训练集上的表现只能说明模型有没有记住训练数据，不能直接说明它是否能泛化。切分方式、交叉验证、类别不平衡和阈值选择，都会改变评估结论。

第三道防线是 \textbf{scientific interpretation}。一个模型可以很好地预测标签，却仍然不能证明因果关系；一个异常分数可以很高，却仍然可能只是坏数据、边界样本或预处理失败。

\begin{warnbox}[机器学习不会替你定义科学问题]
模型可以优化损失函数，但不会自动知道哪个错误更有科学代价、哪个特征可能泄漏答案、哪个样本不代表真实总体。这些判断必须由研究者在工作流中明确写出来。
\end{warnbox}

\section*{本部分的模型报告底线}

从 Part II 进入 Part III 后，报告对象会从“一个数据处理结果”扩展为“一个模型实验结果”。但底层要求没有改变：结果必须带着来源、方法、诊断和限制一起交付。一个模型分数如果离开了数据入口、baseline、切分方式、错误结构和适用边界，就只是一个漂亮数字。

因此，本部分每个模型实验至少要说明五件事：任务到底是什么，数据和标签从哪里来，baseline 是什么，评估设计如何避免自欺，错误样本怎样回到科学解释。只有这五件事同时存在，模型结果才有资格进入 Part IV 的案例项目。

\section*{进入 Part IV 前应该带走什么}

读完 Part III 后，你应该能把一个小型科学数据集整理成机器学习任务，并写清楚：

\begin{itemize}
    \item 要预测或发现的对象是什么；
    \item 特征、标签和样本来自哪里；
    \item baseline 是什么，为什么它是合理起点；
    \item 训练、验证、测试如何划分；
    \item 主要错误类型是什么，哪些错误更有科学代价；
    \item 模型结果在什么范围内可信，在什么范围内应当降级为候选线索。
\end{itemize}

只有做到这些，Part IV 的天文与物理案例才不会变成“套模型”，而会成为真正的端到端小项目。
